% GENERATED FILE. Edit canonical lessons, then run npm run generate:books.
% canonical-source-hash: fnv1a64:6b8bf7ccf16fe8a8
% canonical-lessons: GU-C274-nahi-avu, GU-C274-javano-hato, GU-R274-plans-that-changed

\chapter{What won't happen, and what was going to}
\label{ch:gu-what-won-t-happen-and-what-was-going-to}
\hlchaptermodality{274}

\begin{chapteropening}
\textbf{By the end of this chapter:} \emph{I can say what will not happen (hũ nahī̃ āvũ, te nahī̃ jāy), and what was going to happen but did not (hũ javāno hato, paṇ …).}

Complete GU-R274-plans-that-changed: say what will not happen, and what was going to happen.
\end{chapteropening}

\section[\texorpdfstring{\gu{હું} \gu{નહીં} \gu{આવું} (hũ nahī̃ āvũ)}{હું નહીં આવું (hũ nahī̃ āvũ)}]{\gu{હું} \gu{નહીં} \gu{આવું} (hũ nahī̃ āvũ) — won't}

\label{lesson:GU-C274-nahi-avu}



\begin{quote}
Say \textquotedblleft{}I will come tomorrow\textquotedblright{} and \textquotedblleft{}she will go\textquotedblright{}. (\emph{hũ kāle āvīś}, \emph{te jaśe}.) Now the same two, but they won't happen.
\end{quote}

\begin{grammarlens}[title={Grammar Lens: \gu{નહીં} and a short form}]
\begin{quote}
\textbf{\gu{હું} \gu{કાલે} \gu{નહીં} \gu{આવું}.}
\end{quote}

\begin{quote}
\emph{hũ kāle nahī̃ āvũ.} — \textquotedblleft{}I won't come tomorrow.\textquotedblright{}
\end{quote}

\textquotedblleft{}Won't\textquotedblright{} is \textbf{\gu{નહીં}} with a short form of the verb: no \textbf{-\gu{ીશ}}, no \textbf{-\gu{શે}}. Learn it as one piece. The short form follows the person, and it is the same for a man and a woman:

\noindent
\begin{tabularx}{\linewidth}{@{}>{\raggedright\arraybackslash}X>{\raggedright\arraybackslash}X>{\raggedright\arraybackslash}X@{}}
\toprule
\textbf{who} & \textbf{\gu{આવવું}} & \textbf{\gu{કરવું}} \\
\midrule
\textbf{\gu{હું}} & \textbf{\gu{નહીં} \gu{આવું}} & \textbf{\gu{નહીં} \gu{કરું}} \\
\textbf{\gu{તું}}, \textbf{\gu{તે}}, \textbf{\gu{તેઓ}} & \textbf{\gu{નહીં} \gu{આવે}} & \textbf{\gu{નહીં} \gu{કરે}} \\
\textbf{\gu{અમે}}, \textbf{\gu{આપણે}} & \textbf{\gu{નહીં} \gu{આવીએ}} & \textbf{\gu{નહીં} \gu{કરીએ}} \\
\textbf{\gu{તમે}} & \textbf{\gu{નહીં} \gu{આવો}} & \textbf{\gu{નહીં} \gu{કરો}} \\
\bottomrule
\end{tabularx}
\end{grammarlens}

\subsection*{You'll want to know: the short verbs}
\begin{quote}
\textbf{\gu{તે} \gu{આજે} \gu{નહીં} \gu{જાય}.}
\end{quote}

\begin{quote}
\emph{te āje nahī̃ jāy.} — \textquotedblleft{}She won't go today.\textquotedblright{}
\end{quote}

\textbf{\gu{જવું}} gives \textbf{\gu{નહીં} \gu{જાઉં}} (\emph{nahī̃ jāũ}), \textbf{\gu{નહીં} \gu{જાય}}, \textbf{\gu{નહીં} \gu{જઈએ}}, \textbf{\gu{નહીં} \gu{જાઓ}}. \textbf{\gu{લેવું}} gives \textbf{\gu{નહીં} \gu{લઉં}}, \textbf{\gu{નહીં} \gu{લે}}; \textbf{\gu{ખાવું}} gives \textbf{\gu{નહીં} \gu{ખાઉં}}, \textbf{\gu{નહીં} \gu{ખાય}}. \textbf{\gu{થવું}} gives \textbf{\gu{નહીં} \gu{થાય}}: \textbf{\gu{મોડું} \gu{નહીં} \gu{થાય}} (\emph{moḍũ nahī̃ thāy}), \textquotedblleft{}it won't be late\textquotedblright{}.

Keep these short forms after \textbf{\gu{નહીં}} for now. \textquotedblleft{}I will\textquotedblright{} still takes \textbf{-\gu{ીશ}}: \textbf{\gu{હું} \gu{જઈશ}}, but \textbf{\gu{હું} \gu{નહીં} \gu{જાઉં}}.

\subsection*{Your turn}
Say these aloud:

\begin{itemize}
\raggedright
  \item \textquotedblleft{}We won't eat at home\textquotedblright{} — \emph{ame ghare nahī̃ khāīe}
  \item \textquotedblleft{}I won't take the medicine\textquotedblright{} — \emph{hũ davā nahī̃ laũ}
\end{itemize}

\begin{tcolorbox}[breakable,colback=teal!4,colframe=teal!35!black,title={Before you move on}]
\textquotedblleft{}I won't do it\textquotedblright{}? (\textbf{\gu{હું} \gu{નહીં} \gu{કરું}}.) \textquotedblleft{}He won't come\textquotedblright{}? (\textbf{\gu{તે} \gu{નહીં} \gu{આવે}}.) \textquotedblleft{}You won't go\textquotedblright{}, with respect? (\textbf{\gu{તમે} \gu{નહીં} \gu{જાઓ}}.) And \textquotedblleft{}I will go\textquotedblright{}? (\textbf{\gu{હું} \gu{જઈશ}}: the short form stays with \textbf{\gu{નહીં}}.)
\end{tcolorbox}
\section[\texorpdfstring{\gu{હું} \gu{જવાનો} \gu{હતો} (hũ javāno hato)}{હું જવાનો હતો (hũ javāno hato)}]{\gu{હું} \gu{જવાનો} \gu{હતો} (hũ javāno hato) — was going to}

\label{lesson:GU-C274-javano-hato}



\begin{quote}
Say \textquotedblleft{}I'm going to go tomorrow\textquotedblright{}, as a plan already made. (\emph{hũ kāle javāno chũ}.) Then \textquotedblleft{}I won't go\textquotedblright{}. (\emph{hũ nahī̃ jāũ}.)
\end{quote}

\begin{grammarlens}[title={Grammar Lens: \gu{છું} becomes \gu{હતો}}]
\begin{quote}
\textbf{\gu{હું} \gu{ગઈકાલે} \gu{બજાર} \gu{જવાનો} \gu{હતો}.}
\end{quote}

\begin{quote}
\emph{hũ gaīkāle bajār javāno hato.} — \textquotedblleft{}I was going to go to the market yesterday.\textquotedblright{}
\end{quote}

Put \textbf{\gu{હતો}} where \textbf{\gu{છું}} stood, and the plan moves into the past. Both halves agree, as \textbf{\gu{જવાનો} \gu{છું}} did: a woman says \textbf{\gu{જવાની} \gu{હતી}}, a group \textbf{\gu{જવાના} \gu{હતા}}.

With a verb that has an object, the doer stays plain: \textbf{\gu{હું} \gu{પત્ર} \gu{લખવાનો} \gu{હતો}} (\emph{hũ patra lakhvāno hato}), not \textbf{\gu{મેં}}. It is a plan, not a past of \textbf{\gu{લખવું}}.
\end{grammarlens}

\subsection*{You'll want to know: ..., \gu{પણ} ...}
\begin{quote}
\textbf{\gu{અમે} \gu{જવાના} \gu{હતા}, \gu{પણ} \gu{વરસાદ} \gu{પડ્યો}.}
\end{quote}

\begin{quote}
\emph{ame javānā hatā, paṇ varsād paḍyo.} — \textquotedblleft{}We were going to go, but it rained.\textquotedblright{}
\end{quote}

A plan in the past often did not happen. \textbf{\gu{પણ}} names what got in the way: \textbf{\gu{મીરા} \gu{આવવાની} \gu{હતી}, \gu{પણ} \gu{ટ્રેન} \gu{મોડી} \gu{હતી}} (\emph{Mīrā āvvānī hatī, paṇ ṭren moḍī hatī}), \textquotedblleft{}Mira was going to come, but the train was late.\textquotedblright{}

\subsection*{Your turn}
\begin{itemize}
\raggedright
  \item \emph{Say it:} \textquotedblleft{}Mira was going to come, but the train was late\textquotedblright{} — \emph{Mīrā āvvānī hatī, paṇ ṭren moḍī hatī}
  \item \emph{Contrast:} \emph{hũ javāno chũ} and \emph{hũ javāno hato}: which plan can still happen
\end{itemize}

\begin{tcolorbox}[breakable,colback=teal!4,colframe=teal!35!black,title={Before you move on}]
\textquotedblleft{}I was going to go\textquotedblright{}? (\textbf{\gu{હું} \gu{જવાનો} \gu{હતો}}, or \textbf{\gu{જવાની} \gu{હતી}}.) \textquotedblleft{}We were going to go, but it rained\textquotedblright{}? (\textbf{\gu{અમે} \gu{જવાના} \gu{હતા}, \gu{પણ} \gu{વરસાદ} \gu{પડ્યો}}.) \textquotedblleft{}I was going to write a letter\textquotedblright{}: \textbf{\gu{મેં}} or \textbf{\gu{હું}}? (\textbf{\gu{હું}}.)
\end{tcolorbox}
\section[(review)]{Review — what won't happen, and what was going to}

\label{lesson:GU-R274-plans-that-changed}



\begin{quote}
Say \textquotedblleft{}Arun has gone home\textquotedblright{}. (\emph{Aruṇ ghare gayo che}.) Then \textquotedblleft{}I will come tomorrow\textquotedblright{}. (\emph{hũ kāle āvīś}.)
\end{quote}

\subsection*{Your turn}
Say these aloud:

\begin{itemize}
\raggedright
  \item I won't come tomorrow. She won't go today.
  \item We won't eat at home. He won't take the medicine.
  \item I was going to go to the market, but it rained.
  \item Mira was going to come, but the train was late.
  \item I was going to write a letter, but I came home late.
\end{itemize}

\begin{tcolorbox}[breakable,colback=teal!4,colframe=teal!35!black,title={Before you move on}]
\textquotedblleft{}I won't go\textquotedblright{}? (\textbf{\gu{હું} \gu{નહીં} \gu{જાઉં}}.) Does it change for a woman? (No.) \textquotedblleft{}We were going to go\textquotedblright{}? (\textbf{\gu{અમે} \gu{જવાના} \gu{હતા}}.) What turns \textbf{\gu{જવાનો} \gu{છું}} into a plan that is over? (\textbf{\gu{હતો}} in place of \textbf{\gu{છું}}.)
\end{tcolorbox}
